\documentclass[12pt]{article}
\usepackage[utf8]{inputenc}
\usepackage[T1]{fontenc}
\usepackage{hyperref}
\title{Soga Language Resources for K-12 Curriculum Development}
\author{Compiled by Claude Sonnet 5.5 (claude-sonnet-5-5) for Andr\'as Kornai}
\date{}
\begin{document}
\maketitle

\section{Introduction}
Soga (Lusoga, ISO 639-3 xog) is spoken by roughly two million people, mainly in the Busoga region of Uganda. The search returned little beyond lexicographic scholarship, so this report is short. It lists only what was found in search results.

\section{Primary Educational Resources}
No Lusoga primers or textbooks turned up in the search. Scholarly work argues that the monolingual dictionary is the only Lusoga reference with enough essential content to start elementary teaching. It was written for advanced users and lacks beginner explanations, so a teacher's guide has been proposed.
\begin{itemize}
\item The Monolingual Lusoga Dictionary Faced with Demands from a New User Category (Lexikos): \url{https://lexikos.journals.ac.za/pub/article/view/146}
\end{itemize}

\section{Linguistic and Reference Materials}
\begin{itemize}
\item Compiling the First Monolingual Lusoga Dictionary (Lexikos). It describes how the orthography, grammar and corpus were built for the dictionary \emph{Eiwanika ly'Olusoga}: \url{https://lexikos.journals.ac.za/pub/article/view/434}
\item The same article on AJOL: \url{https://www.ajol.info/index.php/lex/article/view/62723}
\item Wikipedia entry on Soga, which mentions the Lusoga Language Authority (LULA). Parts of it are not fully sourced: \url{https://en.Wikipedia.com/wiki/Lusoga}
\end{itemize}

\section{Religious and Community Resources}
None were found in the search results.

\section{Digital and Audio Resources}
None were found in the search results.

\section{Curriculum Development Context}
Lusoga has a published monolingual dictionary and an orthography designed for it. No current primers or textbooks were found. The Lusoga Language Authority and Uganda's National Curriculum Development Centre are the obvious places to ask about school materials, but this search did not verify either.
\end{document}